% Options for packages loaded elsewhere
\PassOptionsToPackage{unicode}{hyperref}
\PassOptionsToPackage{hyphens}{url}
\documentclass[
]{article}
\usepackage{xcolor}
\usepackage[margin=1in]{geometry}
\usepackage{amsmath,amssymb}
\setcounter{secnumdepth}{-\maxdimen} % remove section numbering
\usepackage{iftex}
\ifPDFTeX
  \usepackage[T1]{fontenc}
  \usepackage[utf8]{inputenc}
  \usepackage{textcomp} % provide euro and other symbols
\else % if luatex or xetex
  \usepackage{unicode-math} % this also loads fontspec
  \defaultfontfeatures{Scale=MatchLowercase}
  \defaultfontfeatures[\rmfamily]{Ligatures=TeX,Scale=1}
\fi
\usepackage{lmodern}
\ifPDFTeX\else
  % xetex/luatex font selection
\fi
% Use upquote if available, for straight quotes in verbatim environments
\IfFileExists{upquote.sty}{\usepackage{upquote}}{}
\IfFileExists{microtype.sty}{% use microtype if available
  \usepackage[]{microtype}
  \UseMicrotypeSet[protrusion]{basicmath} % disable protrusion for tt fonts
}{}
\makeatletter
\@ifundefined{KOMAClassName}{% if non-KOMA class
  \IfFileExists{parskip.sty}{%
    \usepackage{parskip}
  }{% else
    \setlength{\parindent}{0pt}
    \setlength{\parskip}{6pt plus 2pt minus 1pt}}
}{% if KOMA class
  \KOMAoptions{parskip=half}}
\makeatother
\usepackage{longtable,booktabs,array}
\newcounter{none} % for unnumbered tables
\usepackage{calc} % for calculating minipage widths
% Correct order of tables after \paragraph or \subparagraph
\usepackage{etoolbox}
\makeatletter
\patchcmd\longtable{\par}{\if@noskipsec\mbox{}\fi\par}{}{}
\makeatother
% Allow footnotes in longtable head/foot
\IfFileExists{footnotehyper.sty}{\usepackage{footnotehyper}}{\usepackage{footnote}}
\makesavenoteenv{longtable}
\setlength{\emergencystretch}{3em} % prevent overfull lines
\providecommand{\tightlist}{%
  \setlength{\itemsep}{0pt}\setlength{\parskip}{0pt}}
\usepackage{bookmark}
\IfFileExists{xurl.sty}{\usepackage{xurl}}{} % add URL line breaks if available
\urlstyle{same}
\hypersetup{
  hidelinks,
  pdfcreator={LaTeX via pandoc}}

\author{}
\date{}

\begin{document}

\section{\texorpdfstring{A pre-registered, falsifiable
measurement-methodology bridge from \texttt{styxx} primitives to NIST AI
RMF 1.0 Measure-function subcategories
(v0.1)}{A pre-registered, falsifiable measurement-methodology bridge from styxx primitives to NIST AI RMF 1.0 Measure-function subcategories (v0.1)}}\label{a-pre-registered-falsifiable-measurement-methodology-bridge-from-styxx-primitives-to-nist-ai-rmf-1.0-measure-function-subcategories-v0.1}

\textbf{Author:} Alexander Rodabaugh (Fathom Lab) \textbf{Date:}
2026-05-28 \textbf{Substrate:} styxx 7.7.10 in-development; companion to
\texttt{papers/EU\_AI\_ACT\_COMPLIANCE\_2026.md} (EU AI Act bridge v0.1)
and \texttt{papers/PAPER\_recursive\_discipline\_2026\_05\_27.md}
(recursive-discipline methodology paper v7); module
\texttt{styxx.compliance.nist\_ai\_rmf} ships in the same commit
\textbf{Status:} v0.1 minimum-viable mapping. Companion to the EU AI Act
bridge, applying the same methodology to a parallel jurisdictional
regime (United States NIST AI RMF 1.0). Not legal advice; not a
substitute for an organization's own NIST AI RMF implementation.

\begin{center}\rule{0.5\linewidth}{0.5pt}\end{center}

\subsection{Abstract}\label{abstract}

NIST AI 100-1 (the AI Risk Management Framework 1.0, January 2023)
defines four core functions for AI risk management: Govern, Map,
Measure, Manage. The Measure function is its analytical engine --- 22
subcategories that evaluate AI systems against seven trustworthy
characteristics (valid/reliable, safe, secure/resilient,
accountable/transparent, explainable/interpretable, privacy-enhanced,
fair/bias-mitigated). The RMF is voluntary in US federal contexts but is
increasingly referenced by federal procurement, state-level legislation,
and private-sector contracting. This paper introduces
\texttt{styxx.compliance.nist\_ai\_rmf}, a parallel jurisdictional
bridge to the EU AI Act compliance bridge introduced the same day. It
maps styxx primitives to five Measure subcategories (MS-2.3, MS-2.4,
MS-2.5, MS-2.6, MS-2.13) with calibrated metrics, construct ceilings,
commit-level reproducibility receipts, and the same three kill-gates (A1
specific subcategory citations, A2 construct ceilings disclosed, A3
uncovered ≥ covered) enforced by tests. v0.1 explicitly enumerates six
uncovered Measure subcategories (MS-2.7, MS-2.8, MS-2.9 partial,
MS-2.10, MS-2.11, MS-2.12) with alternative-tool references. The bridge
shares dataclasses with the EU AI Act bridge through
\texttt{styxx.compliance.\_common}, signaling that the methodology
pattern transfers across regulatory regimes. Not legal advice;
independent review required for production deployment.

\begin{center}\rule{0.5\linewidth}{0.5pt}\end{center}

\subsection{1. Background}\label{background}

\subsubsection{1.1 NIST AI RMF 1.0 in
context}\label{nist-ai-rmf-1.0-in-context}

NIST AI 100-1, published January 2023, organizes AI risk management into
four functions: \textbf{Govern} (organizational policies), \textbf{Map}
(context establishment, risk identification), \textbf{Measure}
(analytical evaluation), \textbf{Manage} (risk prioritization and
response). The Measure function operationalizes the framework's
analytical claims via 22 subcategories distributed across four
categories (MS-1.x, MS-2.x, MS-3.x, MS-4.x). MS-2.x specifically
addresses the seven trustworthy AI characteristics enumerated in NIST AI
100-1 §4: valid and reliable; safe; secure and resilient; accountable
and transparent; explainable and interpretable; privacy-enhanced; fair
with harmful bias managed.

The RMF is \textbf{voluntary} in US federal contexts. Its adoption is
nevertheless meaningful: federal procurement contracts increasingly
reference it, state legislatures (notably California, New York,
Colorado) cite it in AI governance bills, and major US enterprises adopt
it as a de-facto standard pending future statutory frameworks.

\subsubsection{1.2 Why a parallel bridge}\label{why-a-parallel-bridge}

The EU AI Act bridge (\texttt{papers/EU\_AI\_ACT\_COMPLIANCE\_2026.md})
addressed one jurisdictional regime (Europe, mandatory, August 2026
enforcement deadline). The NIST bridge addresses a parallel regime (US,
voluntary, no fixed deadline) with a structurally different framework:
NIST organizes by \emph{function} (Govern/Map/Measure/Manage) and
\emph{trustworthy characteristics} (seven), where the EU AI Act
organizes by \emph{article} and \emph{Annex III high-risk category}.
Mapping styxx primitives to \emph{both} frameworks demonstrates that the
underlying methodology --- pre-registration, construct ceilings,
commit-level receipts, kill-gates A1/A2/A3 --- generalizes across
regulatory regimes.

The bridge is operationally useful: US organizations adopting NIST AI
RMF for federal procurement or state-level AI governance can cite the
same styxx primitives for Measure subcategories that EU operators cite
for Article 15 sub-paragraphs. Cross-jurisdictional AI deployments
benefit doubly.

\subsubsection{1.3 What this paper IS, and is
NOT}\label{what-this-paper-is-and-is-not}

\textbf{IS:} a measurement methodology bridge from styxx primitives to
NIST AI RMF Measure subcategories, with the same kill-gate discipline as
the EU AI Act bridge. v0.1 ships under MIT alongside the underlying
primitives. Pre-registered, falsifiable, open-source, commit-level
reproducible.

\textbf{IS NOT:} legal advice. A complete NIST AI RMF implementation
(the RMF requires Govern, Map, Manage processes that styxx primitives do
not address). An endorsement by NIST or any standards body. Sufficient
on its own for any contractual or regulatory NIST RMF compliance claim
--- operators must conduct independent legal and technical review.

\begin{center}\rule{0.5\linewidth}{0.5pt}\end{center}

\subsection{2. Methodology}\label{methodology}

The NIST bridge inherits the methodology from the EU AI Act bridge (see
\texttt{papers/EU\_AI\_ACT\_COMPLIANCE\_2026.md} §2 for the full design
principles). Three structural differences:

\begin{enumerate}
\def\labelenumi{\arabic{enumi}.}
\tightlist
\item
  \textbf{Citation format}: clauses are \texttt{MS-N.N} (e.g.,
  \texttt{MS-2.5}), not \texttt{Article\ N.N(x)}.
\item
  \textbf{Coverage scope}: only the Measure function's MS-2.x
  subcategories. The Govern, Map, Manage functions are out of scope by
  design (styxx is a measurement layer, not a governance/management
  framework).
\item
  \textbf{Shared dataclasses}: \texttt{ComplianceMap},
  \texttt{PrimitiveCoverage}, \texttt{UncoveredRequirement} are imported
  from \texttt{styxx.compliance.\_common} and shared with the EU AI Act
  bridge. Both bridges produce the same dataclass instances; only the
  clause-citation format differs.
\end{enumerate}

The same three kill-gates apply: - \textbf{A1 (validity).} Every clause
cites a specific Measure subcategory; regex
\texttt{\^{}MS-\textbackslash{}d+\textbackslash{}.\textbackslash{}d+\$}.
- \textbf{A2 (falsifiability).} Every
\texttt{PrimitiveCoverage.construct\_ceiling} is non-empty and at least
50 characters. - \textbf{A3 (boundary explicitness).}
\texttt{len(uncovered\_requirements())\ \textgreater{}=\ len(coverage\_table())}.
v0.1 ships 6 uncovered vs 5 covered.

A1, A2, A3 are enforced at test time by
\texttt{tests/test\_compliance\_nist\_ai\_rmf.py}.

\begin{center}\rule{0.5\linewidth}{0.5pt}\end{center}

\subsection{3. The coverage table}\label{the-coverage-table}

{\def\LTcaptype{none} % do not increment counter
\begin{longtable}[]{@{}
  >{\raggedright\arraybackslash}p{(\linewidth - 4\tabcolsep) * \real{0.3333}}
  >{\raggedright\arraybackslash}p{(\linewidth - 4\tabcolsep) * \real{0.3333}}
  >{\raggedright\arraybackslash}p{(\linewidth - 4\tabcolsep) * \real{0.3333}}@{}}
\toprule\noalign{}
\begin{minipage}[b]{\linewidth}\raggedright
Measure subcategory
\end{minipage} & \begin{minipage}[b]{\linewidth}\raggedright
Trustworthy characteristic
\end{minipage} & \begin{minipage}[b]{\linewidth}\raggedright
styxx primitives
\end{minipage} \\
\midrule\noalign{}
\endhead
\bottomrule\noalign{}
\endlastfoot
MS-2.3 (performance demonstrated qualitatively/quantitatively, measures
documented) & valid/reliable & \texttt{cognometric\_card},
\texttt{critique\_detector}, \texttt{gauntlet\ +\ preflight} \\
MS-2.4 (production monitoring of system functionality) & valid/reliable
& \texttt{cognometric\_card}, \texttt{agent\_audit} \\
MS-2.5 (valid and reliable; limitations of generalizability documented)
& valid/reliable & \texttt{cognometric\_card},
\texttt{critique\_detector}, \texttt{gauntlet\ +\ preflight} \\
MS-2.6 (safety: real-time monitoring, fail-safe under knowledge-limit
excursion) & safe & \texttt{cognometric\_card},
\texttt{recover\_posture}, \texttt{agent\_audit} \\
MS-2.13 (effectiveness of TEVV metrics and processes evaluated) & TEVV
process effectiveness & \texttt{agent\_audit} \\
\end{longtable}
}

\subsubsection{3.1 Why MS-2.5 is the most direct
fit}\label{why-ms-2.5-is-the-most-direct-fit}

MS-2.5 reads: \emph{``The AI system to be deployed is demonstrated to be
valid and reliable. \textbf{Limitations of the generalizability beyond
the conditions under which the technology was developed are
documented}.''} (Emphasis added.) The bolded clause is the exact
regulatory hook for styxx's construct-ceiling discipline. Every
\texttt{PrimitiveCoverage.construct\_ceiling} field in the registry
documents a published failure mode (e.g., logprob-validity works for
refusal but fails for hallucination --- confident confabulation;
sycophancy FPR 0.30 on restrained-tech responses). These are not buried
caveats; they are first-class structured fields enforced by kill-gate
A2.

A NIST RMF Measure-function implementer documenting MS-2.5 can cite the
styxx primitives + their construct ceilings as direct evidence of
generalizability-limit documentation.

\subsubsection{3.2 Why MS-2.13 is the recursive
subcategory}\label{why-ms-2.13-is-the-recursive-subcategory}

MS-2.13: \emph{``Effectiveness of the employed TEVV metrics and
processes in the measure function are evaluated and documented.''} This
is the meta-subcategory --- it asks whether the measurement methodology
itself is effective.

The styxx project's
\texttt{papers/PAPER\_recursive\_discipline\_2026\_05\_27.md} v7
documents six layers of self-falsification of the project's own
measurement methodology, with pre-registered kill-gates and public
commit-level receipts. \texttt{styxx.agent\_audit} is the specific
instrument that operationalizes this for substrate-checkable claims
(Layer 5 of the arc). The MS-2.13 mapping is unusually clean: the
recursive-discipline paper IS the documentation MS-2.13 asks for.

\begin{center}\rule{0.5\linewidth}{0.5pt}\end{center}

\subsection{4. What this v0.1 does NOT
cover}\label{what-this-v0.1-does-not-cover}

Per kill-gate A3, this section is intentionally at least as long as §3.
Six Measure subcategories are explicitly uncovered:

\subsubsection{4.1 MS-2.7 (secure and
resilient)}\label{ms-2.7-secure-and-resilient}

styxx instruments observe agent cognition, not system security or
resilience under adversarial conditions. Prompt-injection robustness is
\emph{adjacent} (refusal-related instruments may flag some attempts) but
security is not the primary scope.

\textbf{Alternative:} Dedicated LLM security tools (Lakera Guard,
Prompt-Armor, OWASP LLM Top 10 mitigations); standard
application-security audit; penetration testing; red-teaming.

\subsubsection{4.2 MS-2.8 (accountable and
transparent)}\label{ms-2.8-accountable-and-transparent}

styxx produces measurement evidence; it does not generate
deployer-facing accountability documentation, decision-rationale
interfaces, or governance artifacts that MS-2.8 implies.

\textbf{Alternative:} Model card frameworks; capability-statement
templates; deployer documentation processes; AI ethics review boards.

\subsubsection{4.3 MS-2.9 (explainable and interpretable) --- partial
gap}\label{ms-2.9-explainable-and-interpretable-partial-gap}

styxx's construct-ceiling discipline contributes to interpretability by
documenting WHEN each primitive fails, but does not generate
per-decision feature attributions, counterfactual explanations, or
per-input rationales that modern XAI methods provide.

\textbf{Alternative:} Dedicated XAI libraries (SHAP, LIME, Captum,
Inseq); mechanistic interpretability research methods; user studies.

\subsubsection{4.4 MS-2.10
(privacy-enhanced)}\label{ms-2.10-privacy-enhanced}

styxx does not measure differential privacy guarantees, membership
inference resistance, training-data memorization, or PII leakage
detection. Cognometric instruments operate on outputs but not on privacy
properties of those outputs.

\textbf{Alternative:} Differential privacy auditing tools (Google's DP
libraries); PII leakage scanners (Presidio, scrubadub); membership
inference test suites.

\subsubsection{4.5 MS-2.11 (fair with harmful bias
managed)}\label{ms-2.11-fair-with-harmful-bias-managed}

styxx primitives operate on per-step agent cognition signals; they do
not measure population-level disparate impact, demographic outcome
equalization, or training-data bias amplification across protected
classes. Same gap as EU AI Act Article 15.4.

\textbf{Alternative:} Fairness audit libraries (Fairlearn, AIF360);
domain-specific impact assessment; demographic parity / equalized odds
metrics.

\subsubsection{4.6 MS-2.12 (environmental impact and
sustainability)}\label{ms-2.12-environmental-impact-and-sustainability}

styxx instruments operate at inference time on agent outputs; they do
not measure training-energy consumption, inference-carbon footprint, or
supply-chain environmental impact.

\textbf{Alternative:} ML CO2 calculators (CodeCarbon, MLCO2);
cloud-provider sustainability reports; supply-chain LCA tooling.

\begin{center}\rule{0.5\linewidth}{0.5pt}\end{center}

\subsection{5. Pre-registered falsification
criteria}\label{pre-registered-falsification-criteria}

Same five criteria as the EU AI Act bridge
(\texttt{papers/EU\_AI\_ACT\_COMPLIANCE\_2026.md} §5), adapted to the
NIST context:

\begin{itemize}
\tightlist
\item
  \textbf{F1 (registry validity):} if any clause cites a Measure
  subcategory that does not exist in NIST AI 100-1, this paper is
  falsified.
\item
  \textbf{F2 (calibrated-metric reproducibility):} if any cited
  AUC/accuracy number cannot be reproduced from the cited commit hash
  within ±0.01, that \texttt{PrimitiveCoverage} is falsified.
\item
  \textbf{F3 (construct-ceiling discipline):} if a reviewer produces a
  published failure mode of a styxx primitive that is NOT mentioned in
  any \texttt{construct\_ceiling} field, that entry is falsified.
\item
  \textbf{F4 (boundary violation):} if a styxx primitive is later shown
  to produce evidence for an MS subcategory in
  \texttt{UNCOVERED\_MEASURE}, the mapping was scoped too narrowly.
\item
  \textbf{F5 (citation):} ≥1 independent citation (academic, federal
  procurement reference, state-level legislative citation, or commercial
  adoption note) by 2027-02-01 or methodology reassessed.
\end{itemize}

\begin{center}\rule{0.5\linewidth}{0.5pt}\end{center}

\subsection{6. Reproducibility appendix}\label{reproducibility-appendix}

{\def\LTcaptype{none} % do not increment counter
\begin{longtable}[]{@{}
  >{\raggedright\arraybackslash}p{(\linewidth - 4\tabcolsep) * \real{0.3333}}
  >{\raggedright\arraybackslash}p{(\linewidth - 4\tabcolsep) * \real{0.3333}}
  >{\raggedright\arraybackslash}p{(\linewidth - 4\tabcolsep) * \real{0.3333}}@{}}
\toprule\noalign{}
\begin{minipage}[b]{\linewidth}\raggedright
artifact
\end{minipage} & \begin{minipage}[b]{\linewidth}\raggedright
commit
\end{minipage} & \begin{minipage}[b]{\linewidth}\raggedright
path
\end{minipage} \\
\midrule\noalign{}
\endhead
\bottomrule\noalign{}
\endlastfoot
this paper & this commit &
\texttt{papers/NIST\_AI\_RMF\_BRIDGE\_2026.md} \\
NIST bridge module & this commit &
\texttt{styxx/compliance/nist\_ai\_rmf.py} \\
shared dataclasses & this commit &
\texttt{styxx/compliance/\_common.py} \\
NIST tests (15) & this commit &
\texttt{tests/test\_compliance\_nist\_ai\_rmf.py} \\
companion EU AI Act paper & \texttt{f10cab0} &
\texttt{papers/EU\_AI\_ACT\_COMPLIANCE\_2026.md} \\
recursive-discipline paper v7 & \texttt{cf14c83} &
\texttt{papers/PAPER\_recursive\_discipline\_2026\_05\_27.md} \\
MS-2.13 receipt (agent\_audit instrument) & \texttt{3c24b5e} &
\texttt{papers/agent-self-audit/FINDING\_agent\_claim\_audit\_2026\_05\_28.md} \\
MS-2.5 / MS-2.3 receipt (Baseline-019 PASS) & \texttt{17fdd97} &
\texttt{submissions/baseline\_019\_openai\_critique/submission.json} \\
MS-2.5 / MS-2.6 receipt (cognometric instruments) & \texttt{cf14c83} &
\texttt{submissions/\_results/leaderboard.json} \\
\end{longtable}
}

NIST AI 100-1 official source:
\url{https://nvlpubs.nist.gov/nistpubs/ai/nist.ai.100-1.pdf}.
Measure-function subcategory text retrieved 2026-05-28 from
\url{https://airc.nist.gov/airmf-resources/airmf/5-sec-core/}.

\begin{center}\rule{0.5\linewidth}{0.5pt}\end{center}

\subsection{7. Limitations + scope of
applicability}\label{limitations-scope-of-applicability}

\begin{enumerate}
\def\labelenumi{\arabic{enumi}.}
\tightlist
\item
  \textbf{Measure function only.} Govern, Map, Manage are out of scope.
  A complete NIST AI RMF implementation requires all four functions.
\item
  \textbf{5 of 22 Measure subcategories covered.} MS-2.1, MS-2.2 (test
  set documentation, human-subject evaluation) and MS-1.x, MS-3.x,
  MS-4.x are out of scope for v0.1.
\item
  \textbf{Cross-jurisdictional, not cross-framework.} styxx primitives
  that satisfy NIST MS-2.5 may not satisfy EU AI Act Article 15.1(a) ---
  the citation evidence may overlap but the contractual claims differ.
  Operators must conduct framework-specific conformity assessments.
\item
  \textbf{Voluntary regime.} NIST AI RMF is voluntary in US federal
  contexts. The bridge provides evidence; uptake depends on operator's
  chosen compliance posture.
\end{enumerate}

\begin{center}\rule{0.5\linewidth}{0.5pt}\end{center}

\subsection{8. Conclusion}\label{conclusion}

\texttt{styxx.compliance.nist\_ai\_rmf} v0.1 demonstrates that the EU AI
Act bridge methodology transfers to a parallel jurisdictional regime
with a different organizational structure. Both bridges share
dataclasses, share the same kill-gate discipline, and reference the same
underlying styxx primitives. v0.1 covers five MS-2.x Measure
subcategories with five styxx primitives, enumerates six uncovered
subcategories with alternative tool references, enforces structural
integrity at test time, and pre-registers five falsification criteria.

Cross-jurisdictional AI deployments --- increasingly common as US and EU
operators serve the same downstream markets --- benefit from a single
open-source measurement primitive set that maps to both regimes. The
styxx.compliance namespace establishes that pattern.

Not legal advice. Independent NIST AI RMF implementation review required
for any production deployment.

\begin{center}\rule{0.5\linewidth}{0.5pt}\end{center}

\subsection{Acknowledgments}\label{acknowledgments}

Written 2026-05-28 alongside \texttt{styxx.compliance.nist\_ai\_rmf}
v0.1 in the same session as the companion EU AI Act bridge. Subcategory
text retrieved from the NIST AIRC web resource on the same day. The
recursive-discipline methodology of
\texttt{papers/PAPER\_recursive\_discipline\_2026\_05\_27.md} v7 informs
the pre-registration kill-gates and the construct-ceiling discipline.

This paper is offered to NIST AI RMF stakeholder communities, the AI
Safety Institute Consortium (AISIC), federal procurement architects
citing the RMF, state-level AI governance bill authors, and US
enterprise compliance teams as an open-source measurement-methodology
contribution. Comments, falsifications, and improvements are explicitly
invited via the \texttt{fathom-lab/styxx} GitHub issue tracker.

\end{document}
